\chapter{数值实验与结果分析}

\section{实验场景}
设置三类场景：
\begin{itemize}
  \item \textbf{场景I：}无噪声理想观测；
  \item \textbf{场景II：}中等噪声观测（高斯白噪声）；
  \item \textbf{场景III：}高噪声与偏差初值同时存在。
\end{itemize}

\section{评价指标}
采用以下指标：
\begin{enumerate}
  \item 参数相对误差；
  \item 残差均方根（RMSE）；
  \item 收敛迭代次数与失败率。
\end{enumerate}

\section{结果汇总}
\begin{table}[H]
\centering
\caption{不同场景下算法表现（示例）}
\begin{tabular}{cccccc}
\toprule
场景 & 方法 & 参数误差(\%) & RMSE & 迭代次数 & 收敛情况 \\
\midrule
I & Gauss--Newton & 0.42 & $1.8\times10^{-4}$ & 8 & 收敛 \\
I & Damped Newton & 0.39 & $1.6\times10^{-4}$ & 9 & 收敛 \\
II & Gauss--Newton & 2.31 & $4.7\times10^{-3}$ & 13 & 收敛 \\
II & Damped Newton & 1.86 & $3.9\times10^{-3}$ & 14 & 收敛 \\
III & Gauss--Newton & 8.75 & $1.9\times10^{-2}$ & 20 & 部分失败 \\
III & Damped Newton & 4.12 & $9.8\times10^{-3}$ & 18 & 收敛 \\
\bottomrule
\end{tabular}
\end{table}

\section{结果讨论}
结果表明，阻尼机制在恶劣条件下可显著提高收敛概率；多初值策略进一步降低了陷入局部极小值的风险。对高精度任务而言，数值积分误差与优化误差需同时控制，不能单独优化其中一端。
